\section{Introduction}\label{sec:intro}

Service robots are moving into spaces that were designed for people:
classrooms, kitchens, hospital wards and offices. In these settings a robot
that only executes explicit commands is of limited use. People rarely address a
machine with formal instructions; they wave, point, frown, give a thumbs-up or
simply slump at a desk, and they expect the systems around them to respond
sensibly. Humans read such behaviour by combining several channels at once
\citep{argyle1988bodily}: a facial expression qualifies a gesture, and the
setting in which both occur changes what they mean. The same raised hand is a
request for help in a classroom and a greeting in a kitchen; a thumbs-up given
with an angry face is not approval. For a robot to respond appropriately it must
therefore infer the intent behind a behaviour, not merely recognise the
behaviour itself.

Research on human--robot interaction (HRI) has responded by combining several
perceptual cues. Recent systems fuse facial emotion with task context
\citep{mamodiya2025adaptive}, gestures with speech \citep{chen2022real}, or
several sensing modules within a single collaboration framework
\citep{paul2023enhancing}, and surveys of the field describe a clear shift
from single-cue recognisers towards multimodal, perception-driven decision
making \citep{zhao2025multimodal}. These systems share a common design order.
The set of intents a robot should understand is fixed first, from the point of
view of the interaction designer, and the perception models are then trained
to supply the evidence. Perception is treated as an input that is good enough,
and accuracy is usually reported on footage close to the data the models were
trained on.

This paper argues that the order should be reversed. A label space defined
independently of perception can ask the system for distinctions its sensors
cannot reliably make, and when that happens no fusion model, however capable,
can recover. We encountered exactly this failure while building a
camera-based intent recognition system for a service robot. A long series of
model-side improvements to the fusion stage each gained one or two percentage
points and none overtook a hand-written rule cascade. Measuring the cue models
on held-out footage, rather than on the footage they had been fine-tuned on,
explained why. Facial emotion recognition, which appeared 97\,\% accurate on
its own fine-tuning footage, was 76\,\% accurate on the held-out test
recordings, and its most frequent errors fell precisely on boundaries the
intent definitions relied on: disgust read as anger, and neutral read as sad.
Body-motion recognition was correct for only 43.5\,\% of the clips it had not
been fine-tuned on. The intent vocabulary was, in effect, asking the
perception stack for information it did not have.

We therefore propose \emph{perception-aware intent design}: the reliability
of each perceptual cue is measured on held-out data first, and the intent
vocabulary is then defined so that every decision rests on distinctions the
cues can actually make. We apply the approach to a system that reads four cues
from a single ordinary camera --- facial emotion, hand gesture, body motion and
the scene --- and maps them onto nine actionable intents and eleven robot
responses. Intents are keyed on the most reliable cue (gesture), emotion is
grouped into superclasses whose internal confusions no longer change the
outcome, the unreliable motion cue is removed from the decision, and the scene
acts as a modifier. A self-attention fusion model then combines the cues while
excluding any cue that is unobserved, and a safety-gated policy turns the
inferred intent into a robot action, stopping immediately whenever an
emergency becomes likely. The complete pipeline is designed to run on an
embedded edge device so that no video leaves the robot.

On 773 held-out recorded test clips, the fusion model reaches 82.0\,\%
accuracy against 76.5\,\% for a rule cascade operating on the same perceived
cues (McNemar exact $p = 1.7 \times 10^{-7}$). It detects 36 of 38
emergencies (F1 0.818 against 0.565 for the rules), resolves conflicting cues
such as a sarcastic thumbs-up more accurately, degrades more gracefully when a
cue is missing, and outperforms the rules for each of the four participants held out in a
leave-one-actor-out analysis. Just as importantly, a set of controlled
ablations shows where these gains come from. Redesigning the intent vocabulary
around measured reliability contributes 6.7 percentage points on its own,
whereas the choice of fusion architecture does not matter among modern
alternatives: our model is statistically tied with cross-attention and
low-rank fusion. The lesson for designers of interactive robots is that the
interaction vocabulary, not the learning machinery, is the decisive design
decision.

This paper makes four contributions:
\begin{enumerate}
  \item A \emph{perception-aware design method} for robot intent vocabularies,
        which grounds every intent boundary in the measured, held-out
        reliability of the perceptual cues, together with evidence that
        evaluating cue models on their own fine-tuning footage inflates
        their apparent accuracy enough to mislead design decisions.
  \item A \emph{multimodal intent corpus} of 2,869 recorded clips covering 62
        designed interaction scenarios in two environments, performed by ten
        participants and annotated for all four cues, in which test scenarios
        contain cue combinations that are absent from training.
  \item A \emph{robust fusion approach} that combines combination-coverage
        training with selective cue dropout and explicit masking of missing
        cues, evaluated against a rule baseline, three single-cue baselines and
        seven alternative fusion methods under conflicting, missing and
        degraded cues.
  \item An \emph{end-to-end interactive system} with a safety-gated action
        policy, designed for and exported to an embedded edge device.
\end{enumerate}

The remainder of the paper is organised as follows.
Section~\ref{sec:related} reviews related work. Section~\ref{sec:system}
describes the system and its cue models, and Section~\ref{sec:dataset} the
corpus. Section~\ref{sec:taxonomy} presents the perception-aware design of the
intent vocabulary and Section~\ref{sec:fusion} the fusion model.
Sections~\ref{sec:evaluation} and~\ref{sec:results} report the evaluation
design and results, Section~\ref{sec:deployment} the deployment on an edge
device, and Section~\ref{sec:discussion} discusses implications and
limitations before Section~\ref{sec:conclusion} concludes.
